\section{Lean formalization}\label{app:lean}
The Lean 4 library formalizes Lemmas~\ref{f:iteration} and~\ref{f:bell},
conditional forms of the full-replacement bound \eqref{eq:finite-reset}
and the associated lifetime estimates, a conditional abstract form of
Theorem~\ref{f:forgetting}, the overlap step in the proof of
Lemma~\ref{q1:l18}, and the conditional certification of
Theorem~\ref{lean:conditional}. Table~\ref{tab:lean} gives the declarations.
The one-dimensional development formalizes Lemmas~\ref{q1:l1} and
\ref{q1:l1.1}, the arithmetic of Lemma~\ref{q1:l18} with its fidelity
inputs supplied as hypotheses, and selected exponent calculations
from Lemmas~\ref{q1:l11} and~\ref{q1:l12}.
Theorems~\ref{thm:exponential}, \ref{thm:four-dimensional}, \ref{thm:main} and~\ref{q1:main}, Proposition~\ref{f:extension}, the constructions, the noise analysis and the compilation are not formalized. Theorem~\ref{f:forgetting} is checked from the extension onward: its Lean form assumes the two-output extension of Proposition~\ref{f:extension} and the full-replacement decomposition on the small tori as explicit hypotheses.
Theorem~\ref{lean:conditional} gives the conditional certification statement and proof.

\begin{table}[!htbp]
\caption{The Lean statements.}\label{tab:lean}
\small\centering
\begin{tabular}{@{}>{\raggedright\arraybackslash}p{0.37\textwidth}>{\raggedright\arraybackslash}p{0.59\textwidth}@{}}
\toprule
Paper statement & Lean declaration\\
\midrule
Lemma~\ref{f:iteration}, for every positive \(j\)
& \texttt{QuantumScope.iteration\_extension}\\\addlinespace[3pt]
Lemmas~\ref{f:iteration} and~\ref{f:bell}, the bound \(1/2+2^{-j-1}\)
& \texttt{QuantumScope.iterated\_cloning\_bound}\\\addlinespace[3pt]
Lemma~\ref{f:bell}, for every \(k>0\), CPTP encoder and common decoder
& \texttt{QuantumScope.bell\_sharing\_bound}\\\addlinespace[3pt]
Theorem~\ref{f:forgetting}, the conditional abstract form
& \texttt{QuantumScope.uniform\_forgetting}\\\addlinespace[3pt]
Conditional form of \eqref{eq:finite-reset}, for round or layer channels
with the stated replacement decomposition and weight at least \(p\)
& \texttt{QuantumScope.replacement\_ceiling\_pow},
  \texttt{QuantumScope.replacement\_ceiling}\\\addlinespace[3pt]
Conditional layer estimate under the same hypothesis with \(p>0\),
implying \(\tau\le\log(6/5)/p<1/p\) and hence
\eqref{f4:lifetime-upper} and the bound in Section~\ref{q1:effective}
& \texttt{QuantumScope.replacement\_lifetime\_layer}\\\addlinespace[3pt]
The overlap step in the proof of Lemma~\ref{q1:l18}
& \texttt{QuantumScope.good\_event\_fidelity\_union}\\\addlinespace[3pt]
Theorem~\ref{lean:conditional}, the conditional abstract form
of Theorem~\ref{a:complete}
& \texttt{CertificationScope.abstract\_window\_certification}\\
\bottomrule
\end{tabular}
\end{table}

The channel definitions use finite complex matrices and complex-linear
maps. Complete positivity is positivity after amplification by every
finite ancillary dimension. Trace preservation holds for every matrix.
Every partial-trace marginal of the extension equals the given channel
on all inputs. Entanglement fidelity uses the normalized Bell projector
onto \((|00\rangle+|11\rangle)/\sqrt2\) and leaves the reference qubit
untouched. The conditional certification theorem assumes the physical storage and forgetting results as explicit hypotheses. The conditional forgetting theorem likewise assumes the extension and the replacement decomposition; a full replacement is a convex decomposition \(\rho\mapsto q\Tr(\rho)\sigma+(1-q)\Phi(\rho)\) of the layer channel with \(\Phi\) CPTP.

\paragraph{Library and build.}
In the archive, the \texttt{lean/} subdirectory contains \texttt{lean-toolchain}, \texttt{lake-manifest.json} and the \texttt{QuantumStorage} library. \texttt{Challenge.lean} specifies the fourteen theorems, and \texttt{Solution.lean} imports the proved declarations.
The toolchain is Lean \texttt{4.35.0}. The Mathlib commit is
\begin{verbatim}
56a005855ea17e0b0958f73ec801ae140b0a1679
\end{verbatim}
\needspace{6\baselineskip}
With the pinned toolchain installed (elan installs it from \texttt{lean-toolchain}), run these commands from \texttt{lean/} in a POSIX shell. The optional cache step fetches prebuilt dependencies. Without it,
\texttt{lake build} compiles Mathlib from source, which can take hours.
\begin{verbatim}
export LEAN_NUM_THREADS=1
lake exe cache get
lake build
lake env lean -j1 -M5000 Solution.lean
\end{verbatim}
Each of the fourteen axiom sets printed by the last command consists exactly of \texttt{propext}, \texttt{Classical.choice} and \texttt{Quot.sound}. The expected output, wrapped to fit, is
{\small
\begin{verbatim}
'CertificationScope.abstract_window_certification' depends on axioms:
  [propext, Classical.choice, Quot.sound]
'QuantumScope.bell_sharing_bound' depends on axioms:
  [propext, Classical.choice, Quot.sound]
'QuantumScope.equalMarginalExtension_comp' depends on axioms:
  [propext, Classical.choice, Quot.sound]
'QuantumScope.iteration_extension' depends on axioms:
  [propext, Classical.choice, Quot.sound]
'QuantumScope.iterated_cloning_bound' depends on axioms:
  [propext, Classical.choice, Quot.sound]
'QuantumScope.replacement_fidelity' depends on axioms:
  [propext, Classical.choice, Quot.sound]
'QuantumScope.replacement_ceiling' depends on axioms:
  [propext, Classical.choice, Quot.sound]
'QuantumScope.replacement_ceiling_pow' depends on axioms:
  [propext, Classical.choice, Quot.sound]
'QuantumScope.replacement_lifetime' depends on axioms:
  [propext, Classical.choice, Quot.sound]
'QuantumScope.replacement_lifetime_layer' depends on axioms:
  [propext, Classical.choice, Quot.sound]
'QuantumScope.replacement_eventually' depends on axioms:
  [propext, Classical.choice, Quot.sound]
'QuantumScope.uniform_forgetting' depends on axioms:
  [propext, Classical.choice, Quot.sound]
'QuantumScope.good_event_fidelity' depends on axioms:
  [propext, Classical.choice, Quot.sound]
'QuantumScope.good_event_fidelity_union' depends on axioms:
  [propext, Classical.choice, Quot.sound]
\end{verbatim}
}
Any other axiom, including \texttt{sorryAx}, means the check failed.
\texttt{-M5000} sets the Lean memory budget to 5000 MB for that command.

\paragraph{Checking the statements.}
With the same pinned toolchain and Mathlib, append these lines to
separate copies of \texttt{Challenge.lean} and \texttt{Solution.lean}
in the \texttt{lean/} directory.
\begin{verbatim}
#check @QuantumScope.bell_sharing_bound
#check @CertificationScope.abstract_window_certification
#check @QuantumScope.equalMarginalExtension_comp
#check @QuantumScope.iteration_extension
#check @QuantumScope.iterated_cloning_bound
#check @QuantumScope.replacement_fidelity
#check @QuantumScope.replacement_ceiling
#check @QuantumScope.replacement_ceiling_pow
#check @QuantumScope.replacement_lifetime
#check @QuantumScope.replacement_lifetime_layer
#check @QuantumScope.replacement_eventually
#check @QuantumScope.uniform_forgetting
#check @QuantumScope.good_event_fidelity
#check @QuantumScope.good_event_fidelity_union
\end{verbatim}
Run each copy separately with \texttt{lake env lean -j1 -M5000},
followed by the copy's filename. Compare the fourteen printed declaration
types with the statements in \texttt{Challenge.lean} and the paper.
The challenge file has fourteen intentional \texttt{sorry} bodies to expose
the statements. The proof and axiom checks use \texttt{Solution.lean}. The archive's \texttt{lean/README.md} describes a mechanical comparison of the proved statements, with every definition they use, against \texttt{Challenge.lean}.

\paragraph{The one-dimensional sparsity estimate.}
A second development, in the archive's \texttt{lean-1d/} subdirectory,
formalizes Lemma~\ref{q1:l1} and the classical inputs it uses from
\cite{DMNoise1D2026}. It verifies the sparsity definition
\cite[Definition~\ref*{cl:def:sparse}]{DMNoise1D2026} and its basic
properties, the disjoint-child witness construction of
\cite[Lemma~\ref*{cl:lem:sparse-witness}]{DMNoise1D2026} with its finite
tree counts, and the independent-flag probability bounds of that lemma.
For independent elementary keys with possibly unequal rates at most
\(\eps\), finite stencils of at most \(s_R\) keys and chosen integer radius
bounds, it proves the bound \(z_{R,j}(\eps)\) of Lemma~\ref{q1:l1} with the
displayed constants, and \(C_R \eps^{c_R}\le 1/2\) for
\(0\le \eps\le \eps_{R0}\). It also proves Lemma~\ref{q1:l1.1} with integer primitive-stencil
radius bounds: its explicit cutoff \(\eps_R(d,\delta)\) is positive,
is computable as a rational from \(\lceil d\rceil\) and \(\delta\),
and makes the clock moment summable with sum at most \(\delta\).
It proves the recursion and horizon arithmetic of Lemma~\ref{q1:l18}.
Its fidelity calculation assumes the failure probabilities, the union
bound and the overlap inequality of that proof. The overlap inequality
is established in the companion \texttt{lean/} library. The
one-dimensional development also verifies selected identities and
inequalities in the exponent calculations of Lemmas~\ref{q1:l11}
and~\ref{q1:l12}, as listed in \texttt{lean-1d/README.md}.
Choosing integer radius bounds when the final
circuit is enumerated, which flooring the physical coordinates permits
without changing the constants, and identifying the physical noise with
the product-key model are application steps, not formalized results.
The development also verifies the support-independence and
bounded-overlap estimates underlying
\cite[Lemma~\ref*{cl:lem:joint}]{DMNoise1D2026}, the raw-flag
nested-test sum in its first conclusion, and arithmetic parameter
inequalities of \cite{DMNoise1D2026} listed in \texttt{lean-1d/README.md}. It does not formalize
Lemma~\ref{q1:l1.2}, the probability estimates of Lemmas~\ref{q1:l12}
and~\ref{q1:l18}, the other lemmas of
Sections~\ref{q1:model}--\ref{q1:effective}, Theorem~\ref{q1:main}, the
automaton or the physical compilation. Its theorems use only the axioms
\texttt{propext}, \texttt{Classical.choice} and \texttt{Quot.sound}. The toolchain and
Mathlib commit are those above; build and check instructions and a
statement-by-statement comparison with the printed text are in
\texttt{lean-1d/README.md}.
